\section{Verwendete Medien für den Versuch}
\label{ch:Ergebnisse:sec:Material}

%\textcolor{blue}{einleitende Worte}
Für den Versuch wurden verschiedene Medien verwendet.

Für die Befragung der Personen wurde ein Meeting mit Microsoft Teams aufgesetzt, welches aufgezeichnet werden kann. Die Teilnehmer haben einen Fragenkatalog (s. \autoref{ch:Validierung:sec:Material:Nutzerbefragung}) und einen Aufgabenkatalog (s. \autoref{ch:Validierung:sec:Material:Usability_Test}) erhalten. Diese wurden in Form einer Umfrage über die kostenlose Online-Plattform Survio\footnote{Survio. URL: \url{https://www.survio.com/de/} [Zugegriffen am 14.06.2025]} aufbereitet.

Weiterhin wurden für den Versuch die Testpoolmonitoring"=Oberflächen des alten Systems (s. \autoref{ch:IST-Zustand}) sowie die Dashboards des neuen Systems (s. \autoref{ch:Erstellung Dashboard}) genutzt, um diese zu vergleichen. Die Anzeigedaten des neuen Systems wurden hierbei durch zufällige gemockte Daten ersetzt, damit dieselben tagesaktuellen Daten nicht zwei mal abgefragt werden und eine Verfälschung der gemessenen Zeiten verursachen.

% weitere Materialien

%\item Interaktionsprinzipien (Gestaltungskriterien) der ISO 9241-110
%\item Fragebogens nach Prümper (Kurzfassung %\cite{Pruemper:Fragebogen-S:2006}), welcher die Kriterien abdeckt\\
%--> gekürzt auf 10 Fragen, aufgrund eines geringeren Funktionsumfangs\\
%\item Survio - kostenlose Online-Plattform zur Erstellung von Fragebögen
%\item MS Teams - für die Aufzeichnung der Durchführung
%\item altes System - Testpoolmonitoring-Oberflächen \autoref{ch:IST-Zustand}
%\item neues System - mit gemockten Daten, da sonst die selben (tagesaktuellen) Ergebnisse mit beiden Systemen abgefragt worden wären (Verfälschung der Zeitdaten)
%zur Bearbeitung der Aufgaben nicht 

\subsection{\textcolor{MidnightBlue}{Erstellung des Fragebogens für die Nutzerbefragung}}
\label{ch:Validierung:sec:Material:Nutzerbefragung}
Der Fragebogen für die Nutzerbefragung wurde nach dem Vorbild von Jochen Prümper (s. \autoref{ch:SdF:sec:Fragebogen}) erstellt. Das Validierungsverfahren zielt auf die Auswertung der Usability der Dashboard-Lösung im Vergleich zum bestehenden System ab. Da in die ISO-Norm jedoch auch die Auswertung der User-Experience inbegriffen ist, wurde der Fragenkatalog gekürzt. Die Abgrenzung zwischen den Begriffen Usability und User-Experience wird in \autoref{ch:SdF:sec:Usability_vs_User_Experience} dargestellt.

Der Fragenkatalog wurde auf zehn Fragen gekürzt, um den Umfang gering zu halten. Die genauen Begründungen werden in \autoref{Add:Nutzerbefragung_Prümper:Begründung_Filterung} aufgelistet.

Der daraus resultierende Fragebogen ist in \autoref{Add:Nutzerbefragung_Prümper:gekürzter_Fragebogen} aufgeführt und dessen Fragen den Interaktionsprinzipien zugeordnet. Die verwendeten Abkürzungen der Prinzipien sind aus dem Originalfragebogen übernommen und in \autoref{Add:Nutzerbefragung_Prümper:ISO-Abkürzungen} aufgeschlüsselt.

Nachfolgend wird mit der Kurzform 'Frage 1', 'Frage 2' etc. auf den Fragebogen referenziert.

%Auf der Grundlage einer faktorenanalytischen Überprüfung wurde mit dem ISONORM 9241/10-S eine Kurzversion des ISONORM 92410/10 erstellt. Dabei konnten die vorhandenen sieben Skalen bestätigt und ihnen je 3 Items zugeordnet werden. \cite{Pruemper:Fragebogen-S:2006:Paper}

%Der Fragebogen wurde nach dem Vorbild der Nutzerb
%- Vorbild nach Prümper; \cite{Pruemper:Fragebogen-S:2006}\\
%- Begründung warum welche Fragen rausgekürzt wurden, s. \autoref{Add:Nutzerbefragung_Prümper:Begründung_Filterung}\\
%- verwendeter Fragebogen, s. \autoref{Add:Nutzerbefragung_Prümper:gekürzter_Fragebogen}
%-> Aufschlüsselung der Fragen zu den Interaktionsprinzipien der ISO-Norm und die verwendete Reihenfolge\\
%-> nachfolgend wird mit der Kurzform 'Frage 1', 'Frage 2' etc. auf den Fragebogen referenziert

\subsection{Erstellung des Aufgabenkatalogs für den Usability-Test}
\label{ch:Validierung:sec:Material:Usability_Test}

Der Usability-Test wurde auf Grundlage des Fragebogens gestaltet, um alle abgefragten Szenarien abzudecken und Unterschiede zwischen den Systemen für die Bewertung hervorzuheben.

Der daraus resultierende Aufgabenkatalog ist in \autoref{Add:Usability_Test:Aufgabenkatalog} aufgelistet. Er umfasst den täglich benötigten allgemeinen Status der Tests, die Abfrage verschiedener Varianten von absoluten Werten, dessen Ergebnisse jeweils unterschiedlich abzulesen sind, sowie den Wechsel zwischen den beiden Dash\-boards. Ein zusammenfassender Überblick über die Aufgabentypen ist in \autoref{tab:usability_test:Überblick_Aufgabentypen} aufgelistet. Eine detaillierte Auflistung über die Gründe der Aufgabenauswahl und der verwendeten Reihenfolge ist in \autoref{Add:Usability_Test:Begündung_Filterung} aufgeführt.

Nachfolgend wird mit der Kurzform 'Aufgabe 1', 'Aufgabe 2' etc. auf den Aufgabenkatalog des Usability-Tests referenziert.

\begin{table}[]
    \begin{tabular}{|c|c|cc|}
    \hline
    \multirow{2}{*}{Nr} & \multirow{2}{*}{Stichwort} & \multicolumn{2}{c|}{Aufgabentyp}                                                                                                                                                                                           \\ \cline{3-4} 
                        &                            & \multicolumn{1}{c|}{altes System}                                                                                    & neues System                                                                                        \\ \hline
    \multirow{2}{*}{1}  & \multirow{2}{*}{Green Day} & \multicolumn{1}{c|}{\begin{tabular}[c]{@{}c@{}}sind rote Felder mit 'ERR'\\ oder Zahlen != 0 zu sehen?\end{tabular}} & Sind rote Zahlen zu sehen?                                                                          \\ \cline{3-4} 
                        &                            & \multicolumn{2}{c|}{Zusammenfassung von 2 Systemen}                                                                                                                                                                        \\ \hline
    \multirow{2}{*}{2}  & \multirow{2}{*}{Build}     & \multicolumn{1}{c|}{\begin{tabular}[c]{@{}c@{}}Rechenaufgabe\\ (ohne Farbgestaltung)\end{tabular}}                   & \begin{tabular}[c]{@{}c@{}}Ableseaufgabe\\ (rote Zahl von 2 Systemen\\ zusammenzählen)\end{tabular} \\ \cline{3-4} 
                        &                            & \multicolumn{2}{c|}{Zusammenfassung von 2 Systemen}                                                                                                                                                                        \\ \hline
    \multirow{2}{*}{3}  & \multirow{2}{*}{Release}   & \multicolumn{1}{c|}{\begin{tabular}[c]{@{}c@{}}Rechenaufgabe\\ (2 neue Oberflächen)\end{tabular}}                    & \begin{tabular}[c]{@{}c@{}}Ableseaufgabe\\ (2. Dashbard)\end{tabular}                               \\ \cline{3-4} 
                        &                            & \multicolumn{2}{l|}{\begin{tabular}[c]{@{}l@{}}- das System / Oberfläche wird gewechselt\\ - je eine Aufgabe pro System\end{tabular}}                                                                                      \\ \hline
    \multirow{2}{*}{4}  & \multirow{2}{*}{Modi}      & \multicolumn{1}{c|}{Zählaufgabe}                                                                                     & Ableseaufgabe                                                                                       \\ \cline{3-4} 
                        &                            & \multicolumn{2}{l|}{\begin{tabular}[c]{@{}l@{}}- das System / Oberfläche wird zurück gewechselt\\ - je eine Aufgabe pro System\end{tabular}}                                                                               \\ \hline
    \end{tabular}
    \caption{Usability-Test: Überblick über die Aufgabentypen (Zusammenfassung aus \autoref{Add:Usability_Test:Aufgabenkatalog} und \autoref{Add:Usability_Test:Begündung_Filterung})}
    \label{tab:usability_test:Überblick_Aufgabentypen}
\end{table}
